\section*{Introduction}
\label{sec:introduction}

Spatial asset-pricing models organize cross-sectional return dependence through
an interaction matrix and estimate its strength by spatial quasi-likelihood.
The researcher typically supplies $W$ from geography, industry, supply chains,
or news-derived networks and estimates a spatial coefficient $\rho$ conditional
on that choice
\citep{fernandez_spatial_2011,kou_asset_2018,ge_news-implied_2023}.
This approach is powerful once $W$ has been specified, but the spatial equation
begins after the economically relevant relations among firms have already been
chosen.
It therefore leaves two questions outside the model: where the interaction
field comes from, and what $\rho$ measures beyond return dependence conditional
on that field.

Existing constructions encode economic location in different ways.
Point-based spatial models locate a firm in a geographic or characteristic
space, whereas network models represent the firm as a node connected by
observed links.
Text-based finance extends the latter approach by extracting co-coverage and
co-mention relations that balance-sheet classifications can miss
\citep{scherbina_economic_2013,schwenkler_network_2020,ge_news-implied_2023}.
Other work summarizes text as a predictive feature, a named factor, or a learned
graph
\citep{ben_rephael_information_consumption_2019,cong_textual_factors_2024,
son2022graph}.
These approaches establish the relevance of text, but an adjacency still leaves
primitive what constitutes a link and why its cardinal weight should measure
economic interaction.

This paper moves one level upstream of the conventional spatial specification
by changing the mathematical object used to represent a firm.
A diversified firm does not occupy only one economic location: its products,
technologies, supply chains, and information form a footprint across many
positions.
Two firms can share a centroid while having very different footprints, much as
two restaurant chains with the same average store coordinates can cover
different regions.
We therefore represent each firm by a probability measure over
economic-information positions rather than by one point or averaged text score.
A point-valued firm is the special case in which the entire footprint is
concentrated at one position; a genuine distribution also retains spread,
multimodality, and internal composition.

A language model maps each article to a numerical position, so a firm's corpus
forms an empirical distribution in the embedding space.
Once firms are measures, proximity must compare entire footprints rather than
only their centroids.
In the Kantorovich formulation, quadratic optimal transport considers all
feasible couplings of two distributions and selects the one with the smallest
average squared displacement.
Here the ground cost is squared displacement in a language-model embedding
space, so Wasserstein distance measures the least information-space displacement
needed to align one observed footprint with another.
This least-cost formulation gives the geometry an economic interpretation while
making no claim of physical reallocation or literal spatial arbitrage.

Pairwise transport nevertheless does not yet produce an interaction field.
A distance says how far two distributions are separated; it does not say how
several peers jointly represent a fixed target firm.
For each target, optimal transport first aligns its articles separately with
the articles of every candidate peer.
Holding those target-specific alignments fixed, target-anchored Wasserstein
barycentric reconstruction chooses nonnegative unit-sum weights that combine the
aligned peer clouds to approximate the target cloud.
These distributional spanning weights answer which combination of peers jointly
represents the firm, rather than only which single firm lies nearest to it.
Repeating the reconstruction across targets produces $W^\flat$, the directed,
row-stochastic barycentric interaction field.
Its simplex and leave-one-out restrictions give nonnegative unit row sums and a
zero diagonal without a conventional kernel-bandwidth choice.

The paper's primary conceptual contribution is to turn this target-anchored
Wasserstein barycentric reconstruction into an admissible field for spatial
exposure adjustment.
Section~\ref{sec:interaction-operator} formalizes the fixed-target construction
and its boundary with the unrestricted Wasserstein barycentre; the
latter appears
only in Section~\ref{sec:dispersion} as a representation of cross-sectional
dispersion.

Given this field, the paper next asks what the spatial coefficient means.
Each firm has a stand-alone exposure $\xi_i$ implied by its own characteristics,
and a quadratic adjustment problem balances departure from $\xi_i$ against
misalignment with peer exposures.
Solving that problem yields a spatial equilibrium in the peer-adjusted
exposures $B_i$ with feedback coefficient

\[
  \rho=\frac{\lambda}{1+\lambda},
\]

where $\lambda$ is the penalty on peer misalignment relative to the penalty on
departing from the stand-alone exposure.
This spatial closure derives the lag in exposures rather than assuming it in
returns and maps nonnegative adjustment intensity exactly to
$0\leq\rho<1$.
A subsequent return bridge connects the exposure equilibrium to observed
returns.

The same closure permits several admissible fields to enter one adjustment
problem, each with its own coefficient and adjustment index.
The barycentric interaction field and a persistent news co-mention field can
therefore represent separate channels rather than rival estimates of one
privileged network.
A supporting attenuation result then asks how much characteristic-implied
latent exposure dispersion survives peer adjustment under an explicit
maintained transfer restriction.
Related research uses distribution-valued firm characteristics to derive
pairwise covariance restrictions and to study portfolio-risk bounds and
allocation
\citep{gawronsky_continuous_2026,gawronsky_distance_mpt_2026}.
The present paper works at the intervening cross-sectional level: it constructs
the field relating firms and traces how stand-alone exposures propagate through
that field.

Figure~\ref{fig:spatial-exposure-arc} separates what is observed or constructed
from what is latent and model-implied, and from what is estimated.

\begin{figure}[htbp]
  \centering
  \ppfigure{1}{spatial_exposure_arc}
  \caption{From text geometry to model-implied adjustment. Text distributions
    construct the frozen barycentric interaction field $W^\flat$. In the latent
    model, firm-specific information generates stand-alone exposure $\xi_i$,
    which adjustment transforms into peer-adjusted exposure $B_i$. The return
    bridge connects exposure to returns; conditional on the frozen
    field, QMLE yields a
    working-model estimate of $\rho$ and hence the adjustment index
    $\lambda=\rho/(1-\rho)$. Solid arrows denote construction or estimation;
  dashed arrows denote model relations.}
  \label{fig:spatial-exposure-arc}
  \figalttext{Flow diagram with three layers. In the observed and constructed
    layer, firm text distributions produce a frozen peer field; this field and
    returns feed quasi-maximum likelihood estimation. In the latent layer, text
    generates stand-alone exposure, and the peer field adjusts it before a model
    bridge maps exposure to returns. The estimated coefficient is converted to
    an adjustment index. Solid arrows are empirical construction or estimation;
  dashed arrows are maintained model relations, not observed causal paths.}
\end{figure}

In the empirical application, we constructed $W^\flat$ from 2018--2022 text
for \ppvalue{n-tickers-priced} firms and froze it before estimating the return
equation on \ppvalue{post-pooled-2023-2026-w-flat-trading-days} aligned trading
days from 2023 through the incomplete 2026 period.
The pooled QMLE implies the working-model adjustment index
$\hat\lambda=\ppnum[2]{post-pooled-2023-2026-w-flat-lambda}$.
Within the quadratic representation, the fitted penalty on peer misalignment is
therefore \ppnum[2]{post-pooled-2023-2026-w-flat-lambda} times the penalty on
departing from stand-alone exposure.
When the persistent news co-mention field enters jointly, the estimated indices
are $\hat\lambda_B=\ppnum[2]{post-pooled-2023-2026-two-field-lambda-b}$ for the
barycentric field and
$\hat\lambda_N=\ppnum[2]{post-pooled-2023-2026-two-field-lambda-n}$ for the
news-link field.
Each field improves conditional fit once the other is included.

These estimates have a deliberately conditional interpretation.
Freezing the fields before the return window removes mechanical same-sample
feedback, but it does not make the text geometry exogenous to omitted
industries, technologies, attention, or reporting selection.
QMLE estimates working-model return coefficients conditional on the specified
fields and return quasi-likelihood; it does not separately identify the
peer-adjusted exposures or adjustment costs.
Only under the maintained return bridge and quadratic closure does
$\hat\lambda$ inherit the model's adjustment interpretation.
Accordingly, $\hat\lambda$ is a working-model adjustment index, not an observed
managerial cost, a geometry-invariant structural parameter, or a causal peer
effect.

Section~\ref{sec:literature} positions the contribution relative to spatial
finance, network adjustment, text-based measurement, and optimal transport.
Sections~\ref{sec:model} and~\ref{sec:spatial-closure} define stand-alone
exposures and derive the spatial closure, while
Section~\ref{sec:interaction-operator} constructs the barycentric interaction
field and Section~\ref{sec:dispersion} derives the attenuation result.
Section~\ref{sec:identification-estimation} states the return bridge,
identification boundary, data, and estimation design; Section~\ref{sec:results}
reports the evidence, and Sections~\ref{sec:limitations}
and~\ref{sec:conclusion} discuss scope and implications.
